\documentclass[11pt,a4paper]{article}
\usepackage[utf8]{inputenc}
\usepackage[T1]{fontenc}
\usepackage{lmodern}
\usepackage[french]{babel}
\usepackage[a4paper,top=1.5cm,bottom=1.5cm,left=1.8cm,right=1.8cm]{geometry}
\usepackage{parskip}
\usepackage[hidelinks]{hyperref}
\pagestyle{empty}
\setlength{\parindent}{0pt}

\begin{document}

\textbf{Wahid SAMER}\\
Cergy (95), France\\
07 53 10 95 74\\
\href{mailto:samer.ahmed.wahid@gmail.com}{samer.ahmed.wahid@gmail.com}\\
LinkedIn : Wahid SAMER

\vspace{0.4cm}
\begin{flushright}
\textbf{AXA France}\\
Direction Technique Actuarielle Assurance-Vie et Retraite\\
Nanterre\\[0.2cm]
Cergy, le 11 août 2026
\end{flushright}

\vspace{0.2cm}
\textbf{Objet : Candidature -- Chargé d'études actuarielles (F/H)}\\
Référence : AXAFR\_2026030121\_2

\vspace{0.35cm}
Madame, Monsieur,

L'approche ``de bout en bout'' décrite dans votre offre est précisément ce qui m'a conduit à m'intéresser à cette alternance. Comprendre un produit d'épargne ou de retraite, construire les hypothèses qui le décrivent, mesurer sa rentabilité puis étudier l'effet d'une décision sur le risque et la solvabilité constitue, à mes yeux, une application particulièrement concrète des mathématiques financières.

Ma formation est aujourd'hui principalement orientée vers la modélisation quantitative. Elle m'a appris à ne pas considérer un résultat numérique comme une fin en soi : il faut savoir identifier les hypothèses qui le produisent, tester sa sensibilité et comprendre les limites du modèle. Je souhaite désormais appliquer cette démarche à des engagements de long terme, où les comportements clients, les garanties financières et les choix de gestion interagissent dans une même analyse.

C'est également la dimension transversale du poste qui m'intéresse. Le fait de travailler avec les équipes Produit, Contrôle de Gestion et Risk Management oblige à transformer un raisonnement technique en conclusion exploitable pour plusieurs interlocuteurs. Cette exigence correspond à ce que je recherche pour mon M2 : progresser non seulement dans la construction des modèles, mais aussi dans leur interprétation et leur utilisation pour une décision métier.

L'exploitation du Data Lake constitue enfin un axe d'apprentissage qui m'attire particulièrement. Je travaille déjà avec Python pour la modélisation et l'analyse de données ; découvrir PySpark dans un environnement où la donnée sert directement à construire des hypothèses de rentabilité serait une évolution naturelle de ces compétences, sans dissocier la programmation du problème financier traité.

Je serais heureux de rejoindre l'équipe Rentabilité Épargne Retraite Individuelle pour une alternance de douze mois à partir du 21 septembre 2026 et de contribuer aux études de rentabilité, aux travaux de modélisation et à l'analyse des impacts sur les produits d'épargne et de retraite.

Je vous remercie pour l'attention portée à ma candidature et serais ravi de pouvoir échanger avec vous.

Cordialement,

\vspace{0.15cm}
\textbf{Wahid SAMER}

\end{document}
